% Transcription — fonds Alexandre Grothendieck, shelfmark 23, batch 2 (pages 21-40).
% Established from the facsimile archives/batches/23/batch-02.pdf (a mirror of
% https://grothendieck.umontpellier.fr/23.pdf), per the skill
% transcribe-grothendieck. The batch file carries no cover sheet and opens at
% page 21: PDF page N is archivists' page 20+N. Checked against the pencilled
% numbers bottom-left, which read 21, 22, 23, 24 on the first four leaves and
% 28, 30, 31, 32, 34, 37 further on, each at the PDF page of the same number
% less twenty — the alignment holds end to end.
%
% Pass: Opus 5 (claude-opus-5), 2026-09-21 — first pass, unchecked against the pages by a human.
%
% The model is worth naming here because it differs from the rest of the
% folder: batches 1, 3, 4 and 5 were read by Fable 5.1 on the same day, this
% one by Opus 5 at the user's instruction. Nothing else in the file records
% the difference.
%
% What the batch is. Twenty leaves of manuscript notes in his hand, in ink,
% one unrelated typescript leaf, and two otherwise blank leaves carrying only
% a title — those two get no \page{}; their words become section headings,
% each with a \note saying which leaf they come off.
%
%   21-22  The end of the run begun on page 18 of batch 1 (cover leaf 17,
%          « Fibrés des germes d'ordre n de sections, etc. Transformations
%          infinitésimales »): Corollaire 3 to the Proposition of page 20 —
%          the homomorphism S(X̄/X_0) → Is_{X_0}(ξ̄, X̄) is injective — and
%          the reading of X̄/X_0 as a scheme over X_0 with structural sheaf
%          Aut^{X_0}_T(ξ) = G^Λ_{X_0}, ξ̄ being written X_0^{{Λ}}. Page 22
%          carries four lines only, finishing the sentence begun at the foot
%          of page 21.
%          On the join: batch 1's header says page 20 breaks off mid-sentence
%          and runs into page 21. It does not, quite. Page 20 ends on the
%          Corollaire 2 whose last line is cut by the edge of the scan, and
%          page 21 opens a fresh statement, « Corollaire 3 » — the run
%          continues, the sentence does not. The genuine mid-sentence join in
%          this batch is between pages 21 and 22: page 21 ends « Ses notations
%          ... » and page 22 opens « s'interprètent, soit comme un
%          automorphisme de ... », and it is picked up there rather than
%          restarted. Batch 1's file is not edited here; the point is recorded
%          for whoever reconciles the two.
%   23-29  A new run, unannounced by any cover, on jets and principal parts:
%          Δ^{(n)}_{X/S} the n-th infinitesimal neighbourhood of the diagonal
%          seen as an X-scheme by pr_1; π^{(n)}(Y/X/S) = the pullback of
%          X ×_S Y to it; P^{(n)}(Y/X/S) = (δ^n_{X/S})_*(π^{(n)}/Δ^{(n)}),
%          « le fibré des éléments de sections d'ordre n de Y/S sur X/S »;
%          the adjunction Hom_X(X', P^{(n)}) ≃ Hom_{Δ^{(n)}}(X' ×_X Δ^{(n)},
%          π^{(n)}); the two maps d^n and ε^n with ε^n d^n = id; the
%          transition maps φ_{nn'} and Δ^{(∞)}_{X/S}; and on page 28 the
%          reading of Γ(P^n(Y/X/S)/Y) as the connexions d'ordre n of Y/X/S,
%          principal homogeneous under the infinitesimal automorphisms
%          Γ(Y^{{Δ^n_{X/S}}}/Y) when Y/X is smooth. Page 27 restates the
%          construction for rings; page 29 is a scratch leaf cancelled by
%          four long diagonals.
%   30-32  « Transformation infinitésimale de type S' de f » : a couple (u,v)
%          of S'-automorphisms of X_{S'} and Y_{S'} inducing the identity and
%          commuting with f_{S'}; these form a group; its action on sections,
%          w̃(g) = u g_{S'} v^{-1}; the identification Γ(X_{S'}/Y_{S'}) ≃
%          Γ(X_Y^{{Y_{S'}}}/Y); then three « propriétés fonctorielles »
%          numbered 1), 2), 3) — the group action, compatibility with a
%          Y-morphism h, and compatibility with fibre products.
%   33     An otherwise blank leaf inscribed « Exemples, cas particuliers »
%          in his hand, top right: the cover of the run that follows. No
%          \page; its words open the section.
%   34-35  Two counterexamples in characteristic p, written fair and reading
%          nearly end to end. (i) F = aX^p + bY^p − 1, A = k[X,Y]/(F): A is
%          regular and k is algebraically closed in its fraction field, yet
%          A ⊗_k k̄ is not integral. (ii) F = Y² − X^p + a, p ≠ 2: A is
%          regular and its fraction field is separable over k, yet A ⊗_k k̄
%          is not regular, the point Y = 0, X = a^{1/p} being singular.
%          Page 35 finishes the computation of m'/m'² for the second.
%   36     One leaf written in landscape on faintly ruled paper: a summary
%          sheet of criteria — δf = 0 ⟺ f ∈ A[B^p], the iterated form
%          ⟺ f ∈ A[B^{p^k}], and against them Ω¹_{B/A} = 0, flatness, and
%          Δ_{B/A}.
%   37     A typescript leaf, not in his hand: a Bourbaki-style rédaction
%          « n° 215 », its own page « -9 7- », on groupes de trivialité and
%          espèces de structures. Unrelated to the notes in hand and not
%          correspondence; skipped, with no \page.
%   38     An otherwise blank leaf inscribed « Foncteurs différentiels
%          d'ordre n » in his hand, top right. No \page.
%   39-40  The run that leaf covers: a fibred category C over the preschemes
%          on S, a functor F to Φ^C, the objects {X'/S'}, and the example of
%          P^n_{X'/S'}; then on page 40 the schemes S^{d,n} and E^n_{S^{d,n}}
%          = E^n ×_S S^{d,n}. Both leaves are a fast draft, heavily struck,
%          page 40 crossed by two long diagonals; the apparatus there is very
%          dense and honestly so.
%
% Notation fixed for the batch, and continued from batch 1. His underlined
% letters — O_S, O_X, Λ, Γ, S, Aut, Is, Hom — are rendered \underline{...},
% never in prose. The superscripted brace notation X_0^{\{\Lambda\}},
% Y^{\{\Delta^n_{X/S}\}}, X_Y^{\{Y_{S'}\}} is his and is kept as written; it
% names the scheme of sections of type Λ (resp. of type Δ^n, resp. of type
% Y_{S'}) and is introduced at the foot of page 21. He writes the order as a
% parenthesised superscript on Δ and π and as a bare superscript on P, d, ε
% and φ, without ever regularising the two; both are transcribed as they
% stand, and the discrepancy is flagged once, at page 26. « pr_1 », « pr_2 »,
% « Spec », « Aut », « Is », « Hom » are set \operatorname / \mathrm as he
% writes them. The map π^n(Y/X/S) → Y is d^n and the map back on sections is
% ε^n; the two glyphs are near-identical in this hand and the reading is
% settled by the relation ε^n d^n = id on page 25, not by the ink.
%
% Legibility. Pages 30, 34 and 35 are written fair and read almost end to
% end. Pages 21, 23, 24 and 28 are a middling draft whose formulas and
% diagrams carry and whose connective prose often does not. Pages 25, 26, 29,
% 31, 32, 36, 39 and 40 are a fast palimpsest: there the statements come
% through and the sentences between them are largely lost, and the file says
% so rather than inventing them. The two places a checker with the leaves in
% hand will gain most are the boxed aside at the foot of page 30 (three lines
% inside a drawn box, struck by a diagonal) and the top third of page 40
% (three overwritten layers under two diagonals).
%
% What is NOT here. No attempt was made to identify which paragraph of EGA IV
% these notes prepared, nor to match the jets run against EGA IV §16; no date
% appears on the twenty leaves, and none was inferred. The typescript of page
% 37 is described, never transcribed.
\documentclass[11pt,a4paper]{article}
% Le préambule commun à toutes les transcriptions du fonds Grothendieck.
%
% Il tient en une idée : **l'incertitude se compose, elle ne se résout pas.**
% Une transcription qui lisse un mot douteux a détruit la seule chose pour
% laquelle on la faisait — un lecteur doit pouvoir savoir, à chaque endroit,
% ce qui était sur le feuillet et ce qui a été ajouté. Les six macros qui
% suivent sont donc l'appareil critique, et non de la décoration.
%
% Les mêmes commandes sont reconnues par `scripts/render.mjs`, qui produit la
% vue de lecture du site. Une macro ajoutée ici doit l'être aussi là-bas,
% faute de quoi le rendu échouera bruyamment — ce qui est voulu : un rendu
% silencieusement faux est pire qu'un rendu absent.

\ProvidesPackage{grothendieck}[2026/08/08 Transcriptions du fonds Grothendieck]

\usepackage{amsmath,amssymb,amsthm}
\usepackage{mathrsfs}
\usepackage{stmaryrd}
% Les diagrammes commutatifs sont partout dans ce fonds : c'est la langue
% même de la géométrie algébrique de Grothendieck, et une transcription qui
% les rendrait en prose ne transcrirait rien.
\usepackage{tikz-cd}
% Français par défaut — c'est la langue des feuillets — mais l'anglais est
% chargé aussi : la traduction et le résumé sont des documents anglais, et la
% typographie française (espaces avant les deux-points) y serait une faute.
% Ces documents basculent par \selectlanguage{english} en tête de corps.
\usepackage[english,french]{babel}
\usepackage{fontspec}
\usepackage{geometry}
\geometry{a4paper,margin=25mm,marginparwidth=28mm}
\usepackage{marginnote}
\usepackage{xcolor}
% ulem, not soul. `soul`'s \st and \ul are built on a character-by-character
% scan that XeLaTeX's font machinery breaks: the document fails with an
% undefined control sequence rather than degrading. ulem does the same two jobs
% and is Unicode-engine safe. `normalem` stops it hijacking \emph, which the
% transcriptions use for Grothendieck's own emphasis.
\usepackage[normalem]{ulem}
\usepackage{hyperref}
\hypersetup{colorlinks=true,linkcolor=black,urlcolor=black,pdfborder={0 0 0}}

% --- Métadonnées du lot -----------------------------------------------------
% Reprises telles quelles par le script de rendu, qui les affiche en tête de
% la vue de lecture. Elles fixent ce que le fichier prétend couvrir : sans
% elles, une transcription flotte sans point d'attache dans un fonds de seize
% mille feuillets.

\newcommand{\folder}[1]{\gdef\grothFolder{#1}}
\newcommand{\batch}[1]{\gdef\grothBatch{#1}}
\newcommand{\leaves}[2]{\gdef\grothFirst{#1}\gdef\grothLast{#2}}
\newcommand{\foldertitle}[1]{\gdef\grothFoldertitle{#1}}
\gdef\grothFolder{?}\gdef\grothBatch{?}\gdef\grothFirst{?}\gdef\grothLast{?}\gdef\grothFoldertitle{}

% --- Le résumé --------------------------------------------------------------
%
% Il ouvre la lecture modernisée, sous le titre « Résumé », et il est composé
% en petit. Ce n'est pas un ornement : c'est le seul passage du document qui
% ne soit pas des mathématiques, et un lecteur doit pouvoir le reconnaître
% avant de l'avoir lu — puis le sauter s'il connaît déjà le sujet, ou s'y
% arrêter s'il ne le connaît pas. Le filet qui le suit marque la reprise du
% corps.

\newenvironment{resume}
  {\small\setlength{\parskip}{0.45em}}
  {\par\medskip\hrule\medskip}

% --- Mots-clés ---------------------------------------------------------------
%
% En anglais, en clôture du résumé de la lecture modernisée : le vocabulaire
% moderne sous lequel chercher ce que les feuillets construisent. Le manifeste
% du site (`npm run manifest`) les extrait pour étiqueter la cote entière —
% cette ligne est la source unique des tags, il n'y a pas de fichier à côté.
\newcommand{\keywords}[1]{\par\smallskip\noindent
  {\footnotesize\textsc{Keywords} --- \textit{#1}}\par}

% --- L'appareil critique ----------------------------------------------------

% Le numéro de feuillet, en marge. C'est l'ancre qui permet de régler un
% désaccord sur une formule en regardant *un* feuillet plutôt qu'une chemise
% de six cents. Le site s'en sert aussi pour tourner les pages du fac-similé
% quand on fait défiler la transcription.
\newcommand{\leaf}[1]{%
  \marginnote{\footnotesize\color{gray!70}\textsf{#1}}%
  \label{leaf:#1}\ignorespaces}

% Illisible : on ne devine pas. Un mot inventé qui se lit comme les autres est
% le pire résultat possible.
\newcommand{\ill}{\textcolor{red!60!black}{[\,\ldots\,]}}

% Lecture douteuse : le mot est donné, et signalé comme incertain.
\newcommand{\uncertain}[1]{\uline{#1}}

% Ajout de l'éditeur — une lettre manquante, un mot sous-entendu.
\newcommand{\add}[1]{\textcolor{blue!50!black}{[#1]}}

% Biffé par Grothendieck lui-même. Ce qu'il a rayé fait partie du document :
% une rature dit souvent où la pensée a changé de direction.
\newcommand{\struck}[1]{\sout{#1}}

% Note du transcripteur, sur l'état matériel du feuillet.
\newcommand{\note}[1]{\par\smallskip{\small\color{gray!60!black}\textit{[#1]}}\par\smallskip}

% Note marginale de Grothendieck, distincte du corps du texte.
\newcommand{\marginal}[1]{\par\smallskip\noindent\hspace{1em}%
  {\small\color{gray!60!black}#1}\par\smallskip}

% --- Titre ------------------------------------------------------------------

\AtBeginDocument{%
  \begin{center}
    {\large\bfseries Fonds Alexandre Grothendieck --- cote n\textsuperscript{o}~\grothFolder}\\[2pt]
    {\itshape\grothFoldertitle}\\[4pt]
    {\small feuillets \grothFirst--\grothLast\ (lot \grothBatch)}\\[2pt]
    {\footnotesize\color{gray!60!black}Transcription établie d'après le fac-similé numérisé
     par l'Université de Montpellier.\\
     Les lectures incertaines sont soulignées, les illisibles notées [\,\ldots\,],
     les ajouts entre crochets.}
  \end{center}
  \vspace{1em}}

\endinput


\folder{23}
\batch{2}
\pages{21}{40}
\dating{1966-1968}
\watermark{Édition de démonstration}
\foldertitle{EGA IV. Compléments : notes manuscrites (s.d.), lettres (1966, 1968), tapuscrit (s.d.).}

\begin{document}

\section{Fibrés des germes d'ordre $n$ de sections, etc. Transformations
infinitésimales (suite)}
\note{le titre est celui de la chemise vierge qui porte le numéro 17, déjà
donné au lot 1 ; il est répété ici parce que la suite de notes qu'il couvre
enjambe la limite des deux lots. Les pages 21 et 22 achèvent la série de
corollaires commencée page 20}

\page{21}
\textbf{Corollaire 3} On a un hom.\ naturel de faisceaux
\[
  \underline{S}(\overline{X}/X_{0}) \longrightarrow
  \underline{\mathrm{Is}}_{X_{0}}(\overline{\xi}, \overline{X})
\]
\note{le $\overline{\xi}$ est écrit deux fois, la seconde par-dessus la
première et d'un trait plus gras}
compatible avec les opérations de $\underline{\mathrm{Aut}}^{X_0}_{T}(X)$
à gauche, et les opérations \uncertain{à d.} de
$\underline{\mathrm{Aut}}^{X_0}_{T}(\xi)$ [qui opèrent sur le premier
\ill{} par le corollaire 2]. Cet homomorphisme est \emph{injectif} [car
$\underline{\mathrm{Aut}}^{X_0}_{T}(\xi)$ opère fidèlement dans
$\overline{\xi}$]

On peut dire aussi [mais \ill{} n'est pas très utile] que
$\overline{X}/X_{0}$ est muni d'une structure de \struck{fais} schéma sur
$X_{0}$ à faisceau structural :
\[
  \boxed{\;\underline{\mathrm{Aut}}^{X_0}_{T}(\xi) =
  \underline{G}^{\Lambda}_{X_{0}}\;}
\]
opérant sur le \uncertain{schéma-type}
\[
  \xi = X_{0} \times_{S} T = X_{0} \otimes_{\underline{\mathcal{O}}_{S}}
  \underline{\Lambda}
\]

La donnée de $X$ \ill{} sur $T$, se réduisant suivant $X_{0}$, est
équivalente à la donnée d'un schéma $\overline{X}$ sur $X_{0}$, muni d'un
tel faisceau structural.
\note{l'adjectif qui qualifie $X$ sur $T$ est souligné d'un trait et ne se
laisse pas décider sur le film : quatre ou cinq lettres, dont une haute à
l'avant-dernière place}

Le faisceau principal associé n'est autre que
$\underline{S}(\overline{X}/X_{0})$.

$\overline{\xi}$ est noté ainsi $X_{0}^{\{\Lambda\}}$, et s'appelle \ill{}
\struck{\ill{}} des pts propres de $X_{0}$ de type $\Lambda$. Ses notations
\struck{\ill{}} \ill{}
\note{les deux dernières lignes du feuillet sont entamées par le bord du
scan ; un mot biffé ouvre la seconde}

\page{22}
s'interprètent, \emph{soit} comme un automorphisme de
\[
  X_{0} \times_{S} T = X_{0} \otimes_{\underline{\mathcal{O}}_{S}}
  \underline{\Lambda}
\]
induisant l'identité sur $X_{0}$, \emph{soit} comme la donnée d'un champ de
sections \struck{\ill{}} \ldots
\note{le feuillet ne porte que ces quatre lignes, qui achèvent la phrase
laissée en suspens au bas de la page 21 ; le reste est vierge. Le signe
$\times$ du premier membre porte une rature ; le passage biffé de la
dernière ligne contient un $\Lambda$}

\subsection{Éléments de sections d'ordre $n$ de $Y/S$ sur $X/S$}
\note{la formule est de lui : elle est écrite en toutes lettres au haut de
la page 23 et reprise, soulignée, dans la note marginale de la page 24.
Aucune chemise ne couvre cette suite, qui commence sans transition au
feuillet 23}

\page{23}

\begin{tikzcd}[column sep=small, row sep=small]
  Y \arrow[d, "p"'] & \pi^{(n)}(Y/X/S) \arrow[d] \\
  S & X \arrow[l, "f"'] & \Delta^{(n)}_{X/S} \arrow[l, "\delta^{n}_{X/S}"']
\end{tikzcd}

On veut définir en « fibrés »
\[
  P^{n}(Y\,/\,X\,/\,S)
\]
des \struck{éléments} « éléments de sections d'ordre $n$ de $Y/S$ sur
$X/S$ ».

($X \times_{S} X$ est muni des deux projections $pr_{1}$, $pr_{2}$ et de
l'immersion diagonale $\Delta_{X}$. On considère \struck{$\delta$}
$\Delta^{(n)}_{X/S}$ le \ill{} voisinage infinitésimal de l'immersion
$\Delta_{X}$ comme un $X$-préschéma \uncertain{grâce à} $pr_{1}$.
$i^{(n)}_{X}$ désigne l'immersion de $\Delta^{(n)}_{X/S}$ \struck{dans}
$X \times_{S} X$.) On considère alors le $\Delta^{(n)}_{X/S}$-schéma
\[
  (pr_{2}\, i^{(n)}_{X})^{*}(Y/X) = \struck{P^{(n)}(Y/X/S)}\;
  \Delta^{(n)}_{X/S} \times_{X \times_{S} X} (X \times_{S} Y)
  = \struck{\ill{}}\; \pi^{(n)}(Y/X/S)
\]
\note{le membre biffé $P^{(n)}(Y/X/S)$ est écrit au-dessus du produit fibré
et relié à lui par une accolade ; c'est la notation qui sera réservée, page
24, à l'image directe}

\[
  X \times_{\Delta^{(n)}_{X/S}} \pi^{n}(Y/X/S) = Y
\]
d'où $Y \to \pi^{n}(Y/X/S)$ \ill{} $X \xrightarrow{\;\Delta_{X}\;}
\Delta^{(n)}_{X/S}$

\begin{tikzcd}[column sep=small, row sep=small, nodes={font=\scriptsize}]
  & X \times_{S} Y \arrow[d] & \pi^{(n)}(Y/X/S) \arrow[l] \arrow[d]
    & Y \arrow[l, "\varepsilon^{n}(Y/X/S)"'] \arrow[d] \\
  X & X \times_{S} X \arrow[l, "pr_{1}"']
    & \Delta^{(n)}_{X/S} \arrow[l, "i^{(n)}_{X/S}"'] \arrow[d]
    & X \arrow[l, "\Delta_{X/S}"'] \arrow[d] \\
  X' & X' \times_{S} X \arrow[l] & (\Gamma_{\varphi})^{(n)} \arrow[l]
    & X' \arrow[l]
\end{tikzcd}
\note{diagramme redessiné d'un croquis : les trois colonnes de droite
descendent par des traits verticaux, la ligne du bas est ajoutée en plus
petit et se rattache à celle du dessus par deux traits seulement}

\[
  d^{(n)}_{Y/X/S} : \pi^{n}(Y/X/S) \longrightarrow X \times_{S} Y
  \longrightarrow Y
\]
et \ill{} \ill{} commutatif \ldots

\begin{tikzcd}[column sep=small, row sep=small]
  Y \arrow[r, "\varepsilon^{n}(Y/X/S)"] \arrow[d] & \pi^{n}(Y/X/S)
    \arrow[r, "d^{n}(Y/X/S)"] \arrow[d] & Y \arrow[d] \\
  X \arrow[r, "\Delta_{X/S}"] & \Delta^{(n)}_{X/S}
    \arrow[r, "\delta^{(n)}_{X/S}"] & X
\end{tikzcd}
\note{le carré est encadré d'un trait à la main ; le sommet inférieur droit
porte un mot biffé avant le $X$}

\page{24}
Ce n'est donc pas un \struck{$X$}-morphisme en général.

Supposons maintenant que $\Delta^{(n)}_{X/S}$ soit \struck{plat, de
présentation} \struck{finie} loc.\ libre sur $X$. Alors on peut former
\[
  \boxed{\;(\delta^{n}_{X/S})_{*}\bigl(\pi^{n}(Y/X/S) /
  \Delta^{(n)}_{X/S}\bigr) = P^{(n)}(Y/X/S).\;}
\]

\marginal{fibré des éléments de sections d'ordre $n$ $Y/S$ sur $X/S$}
\note{note de sa main dans la marge de gauche, écrite en biais de bas en
haut sur trois lignes, chacune soulignée ; elle est reliée par un trait à la
formule encadrée}

Par définition, on aura donc
\[
  \underline{\Gamma}\bigl(\pi^{n}(Y/X/S)/\Delta^{(n)}_{X/S}\bigr) \simeq
  \underline{\Gamma}\bigl(P^{(n)}(Y/X/S)/X\bigr)
\]
et \struck{de façon} plus précisément, si $\varphi : X' \to X$ est un
préschéma sur $X$
\[
  \operatorname{Hom}_{X}\bigl(X', P^{(n)}(Y/X/S)\bigr) \simeq
  \operatorname{Hom}_{\Delta^{(n)}_{X/S}}\bigl(X' \times_{X}
  \Delta^{n}_{X/S},\ \pi^{(n)}(Y/X/S)\bigr)
\]

Un élément de \uncertain{cet ensemble} est donc \uncertain{obtenu} en
\uncertain{prenant} l'immersion
\[
  \Gamma_{\varphi} = \struck{\ill{}}\ \varphi \times_{S} \mathrm{id}_{X} :
  X' \longrightarrow X' \times_{S} X
\]
\note{le membre médian est récrit par-dessus un premier essai ; tel quel,
$\varphi \times_{S} \mathrm{id}_{X}$ ne va pas de $X'$ dans $X' \times_S X$
--- le graphe de $\varphi$ est $(\mathrm{id}_{X'}, \varphi)$. La page n'est
pas corrigée}

en prenant le \uncertain{$n$-ième voisinage} \ill{}
\[
  \Gamma_{\varphi}^{(n)} \longrightarrow X' \times_{S} X
\]
et une section de $X' \times_{S} Y$ \uncertain{là-dessus}

\struck{On va alors} \uncertain{On} a un $X$-morphisme canonique \ill{}
\uncertain{noté} $D^{n}(Y/X/S)$ :

\page{25}

\begin{tikzcd}[column sep=small, row sep=small]
  P^{(n)}(Y/X/S) \arrow[r, "\varepsilon^{n}(Y/X/S)"] \arrow[dr] & Y \arrow[d] \\
  & X
\end{tikzcd}

D'autre part, on a \ill{} \ill{} un \emph{hom.} \struck{application} de
\emph{faisceaux sur $X$}
\[
  d^{n}(Y/X/S) : \underline{\Gamma}(Y/X) \longrightarrow
  \underline{\Gamma}\bigl(P^{(n)}(Y/X/S)/X\bigr)
\]
[tel que \ill{} composé avec
\[
  \varepsilon^{n}(Y/X/S) : \underline{\Gamma}\bigl(P^{(n)}(Y/X/S)/X\bigr)
  \longrightarrow \underline{\Gamma}(Y/X)
\]
soit l'identité : $\varepsilon^{n} d^{n} = \mathrm{id}$]
\note{les deux glyphes $d$ et $\varepsilon$ se ressemblent beaucoup dans
cette main ; c'est la relation $\varepsilon^{n} d^{n} = \mathrm{id}$, et le
sens des deux flèches, qui décident la lecture, non le trait}

\struck{Je fais donc} Par \emph{ces} \ill{} \struck{\ill{}} qui
\uncertain{revient} \ill{} de \uncertain{définir}
\[
  \underline{\Gamma}(Y/X) \longrightarrow
  \underline{\Gamma}\bigl(P^{n}(Y/X/S)/X\bigr)
\]
et de \uncertain{vérifier} \ill{} \ill{} sections locales sur $X$. On
\uncertain{a} \uncertain{un}
\[
  \underline{\Gamma}(Y/X) \longrightarrow \struck{\operatorname{Hom}}\;
  \underline{\Gamma}\bigl(Y \times_{X} \Delta^{(n)}_{X/S} \,/\,
  \Delta^{(n)}_{X/S}\bigr)
\]
On
\[
  \operatorname{Hom}_{X}\bigl(\Delta^{n}_{X/S}, Y\bigr) =
  \underline{\Gamma}\bigl(Y \times_{X} \Delta^{(n)}_{X/S} \,/\,
  \Delta^{(n)}_{X/S}\bigr)
\]
\ill{} sections \uncertain{au-dessus} de $X \times_{S} Y$ \ill{}
$X \times_{S} X$ \ill{} de l'immersion $X \longrightarrow X \times_{S} X$
et les \uncertain{projections}, \ill{} \ill{}
\note{le bas du feuillet est un brouillon rapide : les formules portent,
les phrases qui les relient sont perdues pour l'essentiel}

\page{26}
\struck{$P^{n}_{A/\Lambda}(M)$}
\note{formule isolée en haut à gauche, biffée de deux traits obliques, avec
un $d^{n}$ au-dessus}

\ill{} en \ill{} \ill{} de éléments de sections d'ordre \ill{} sur
$X \times_{S} X$ le \ill{}

Soit $m \geq n$, alors \ill{} \uncertain{biunivoque} avec les \ill{}
d'ordre \ill{} de $X \times_{S} Y$ \ill{} le \ill{} immersion.
\[
  \Delta^{(m)}_{X/S} \longrightarrow \Delta^{(n)}_{X/S}
\]
Les $\Delta^{(n)}_{X/S}$ forment donc un \emph{système inductif} de schémas
sur $X$, dont la limite est \ill{} \ill{} sur $X$, noté
$\Delta^{(\infty)}_{X/S}$.

On a \uncertain{de même}
\[
  \pi^{(n)}(Y/X/S) = \pi^{n'}(Y/X/S) \times_{\Delta^{(n)}_{X/S}}
  \Delta^{(m)}_{X/S}
\]
d'où aussitôt un $X$-\uncertain{homomorphisme} \uncertain{naturel}
\[
  P^{n'}(Y/X/S) \xrightarrow{\;\varphi_{nn'}\;} P^{n}(Y/X/S)
\]
C'est un système \uncertain{transitif} de $X$-\uncertain{homomorphismes},
d'où
\[
  \begin{cases}
    \varepsilon\, \varphi_{on} = \varepsilon^{n} \\
    \varepsilon^{(n)}(Y/X/S) \circ \varphi_{nn'} = \varepsilon^{(n')} \\
    \varphi_{nn'}\, d^{(n')} = d^{(n)}.
  \end{cases}
\]
\note{la première ligne du système est soulignée d'un trait et la deuxième
d'un double trait ; elles disent la même chose et la seconde paraît récrire
la première. L'ordre est noté $(n)$ entre parenthèses sur $\Delta$ et $\pi$,
et $n$ tout court sur $P$, $d$, $\varepsilon$ et $\varphi$ : la page ne
régularise nulle part les deux écritures, et la transcription ne les
régularise pas non plus}

\page{27}
Les sections de $P^{n}(Y/X/S)$ sur $Y$ \ill{} des
$X$-\uncertain{morphismes}, \ill{}

\begin{tikzcd}[column sep=small, row sep=small, nodes={font=\scriptsize}]
  Y^{\{\Delta^{(n)}_{X/S}\}} \arrow[dr, dashed] & & & \\
  P^{n}(Y/X/S) \arrow[r] \arrow[d, "p"'] & Y \arrow[d]
    & \pi^{(n)}(Y/X/S) \arrow[l] \arrow[d] & Y \arrow[l] \arrow[d] \\
  X & X \times_{S} X \arrow[l, "pr_{1}"']
    & \Delta^{(n)}_{X/S} \arrow[l] & X \arrow[l]
\end{tikzcd}
\note{la flèche en pointillé porte l'exposant à accolade
$Y^{\{\Delta^{(n)}_{X/S}\}}$, de la même famille que le
$X_0^{\{\Lambda\}}$ de la page 21 ; la colonne de droite est reliée par un
trait sans pointe que le diagramme rend par un trait double}

\begin{tikzcd}[column sep=small, row sep=small]
  A \arrow[r] & A \otimes A \\
  B \arrow[u] & A' \arrow[u]
\end{tikzcd}
\note{traduction affine de la construction : $A \to A \otimes A$ est tracé
sur la page avec deux flèches parallèles --- les deux structures de
$A$-algèbre --- rendues ici par une seule ; $B$ est une $A$-algèbre et $A'$
le quotient de $A \otimes A$ qui définit $\Delta^{(n)}$}

\begin{tikzcd}[column sep=small, row sep=small, nodes={font=\scriptsize}]
  (B \otimes_{A} A') \otimes A \arrow[d] & A' \otimes A \arrow[l] \arrow[r]
    & A' & A' \otimes B \arrow[l] \\
  (B \otimes_{A} B) \otimes A & B \otimes A \arrow[l] \arrow[r] & B
    & B \otimes B \arrow[l]
\end{tikzcd}
\note{le bas du feuillet reprend en algèbre le diagramme du haut ; il est
tracé rapidement et plusieurs flèches sont sans pointe sur la page, rendues
ici avec pointe faute de pouvoir décider leur sens --- le point est à
vérifier sur le feuillet}

\page{28}
\[
  m \longrightarrow 1 \otimes m
\]

\begin{tikzcd}[column sep=small, row sep=small]
  & Y^{\Delta^{n}_{X/S}} \arrow[dr] & \\
  P^{n}(Y/X/S) \arrow[r] & Y \arrow[d] & \pi^{n}(Y/X/S) \arrow[l] \arrow[d] \\
  & X & \Delta^{n}_{X/S} \arrow[l]
\end{tikzcd}

\[
  \underline{\Gamma}\bigl(Y^{\{\Delta^{n}_{X/S}\}}/Y\bigr)
\]
\uncertain{groupe} des automorphismes \ill{} infinitésimaux d'ordre $n$ des
fibres $Y/X$

\[
  \underline{\Gamma}\bigl(P^{n}(Y/X/S)/Y\bigr)
\]
$=$ \uncertain{groupe} des automorphismes de
$\pi^{n}(Y/X/S)/\Delta^{n}_{X/S}$ qui induisent l'id.\ sur $Y$

$=$ \emph{connexions d'ordre $n$} \ill{} des fibres $Y/X/S$
\note{le second membre est souligné d'un trait gras : c'est la définition
que la page vise}

Si $Y/X$ est \uncertain{lisse}, \ill{}
$\underline{\Gamma}\bigl(P^{n}(Y/X/S)/Y\bigr)$ \ill{} \uncertain{est}
principal homogène \ill{}
$\underline{\Gamma}\bigl(Y^{\{\Delta^{n}_{X/S}\}}/Y\bigr)$, \ill{}
\note{trois lignes rendues ici de façon très lacunaire : la phrase affirme
que les connexions d'ordre $n$ forment un espace principal homogène sous les
automorphismes infinitésimaux d'ordre $n$, mais les mots qui l'articulent ne
se laissent pas tous lire}

\begin{tikzcd}[column sep=small, row sep=small]
  P^{n}(Y/S) & P^{n}(X/S) \arrow[l] \\
  \mathcal{O}_{Y} \arrow[u] & \mathcal{O}_{X} \arrow[l] \arrow[u]
\end{tikzcd}

hom.\ de $\underline{\mathcal{O}}_{Y}$-algèbres augmentées
\[
  P^{n}(Y/X) \longrightarrow f^{*}\bigl(P^{n}(X/S)\bigr)
\]

\begin{tikzcd}[column sep=small, row sep=small]
  f^{*}P^{n}(Y/S) \arrow[d] & P^{n}(Y/S) \arrow[l, "?"']
    \arrow[d] & f^{*}P^{n}(X/S) \arrow[l, Rightarrow] \arrow[d] \\
  \underline{\mathcal{O}}_{X} & \underline{\mathcal{O}}_{X} \arrow[l]
    & \underline{\mathcal{O}}_{X} \arrow[l]
\end{tikzcd}
\note{la flèche du haut au milieu porte un « ? » et, sous elle, le mot
« comment ? » de sa main : la page pose la question et ne la résout pas}

hom.\ de \struck{de} $f^{*}(P^{n}(X/S))$-algèbres
$\underline{\mathcal{O}}_{X}$-augmentées
\[
  P^{n}(Y/S) \longleftarrow f^{*}P^{n}(X/S)
\]
(pas nécessairement augmentées)

\page{29}
\note{le feuillet entier est barré de quatre longs traits diagonaux ; ce
qui se laisse lire en est transcrit, dans l'ordre où cela vient}

\[
  \varphi\bigl(g\, d^{n}f \otimes m\bigr) = m \otimes f\, d^{n}g
\]
\note{sous « $d^{n}f$ » une accolade descend vers « $g \otimes f$ », en
interligne}

\[
  \varphi\bigl(g\, d^{n}(f\,\ill{})\bigr) = f m \otimes \ill{}
\]
\[
  \ill{} \otimes_{A} (A \otimes_{\Lambda} A)
\]
\[
  M \otimes_{\Lambda} A
\]
\[
  \ill{}(A \otimes_{\Lambda} M) \qquad
  \underline{A} \otimes_{\Lambda} S_{\Lambda}(M) \qquad
  \struck{A} \otimes_{\Lambda} S_{A}(M)
\]
\[
  \struck{\operatorname{Hom}_{B}(M, B)}
\]

\begin{tikzcd}[column sep=small, row sep=small]
  M \arrow[d] \\
  A \arrow[r] & B
\end{tikzcd}
\note{les deux dernières formules sont dans le coin inférieur droit ; le
$\operatorname{Hom}_{B}(M,B)$ est entièrement raturé}

\subsection{Transformations infinitésimales de type $S'$ d'un morphisme}
\note{la formule est de lui : elle ouvre la page 30, soulignée d'un trait.
Aucune chemise ne sépare cette suite de la précédente}

\page{30}
Soit $f : X \longrightarrow Y$ un $S$-morphisme. On appelle
\emph{transformation infinitésimale de type $S'$ de $f$} un couple $u, v$ de
$S'$-automorphismes de $X_{S'}$, $Y_{S'}$ (induisant l'identité sur $X$,
$Y$) rendant commutatif le diagramme

\begin{tikzcd}
  X_{S'} \arrow[r, "u"] \arrow[d, "f_{S'}"'] & X_{S'} \arrow[d, "f_{S'}"] \\
  Y_{S'} \arrow[r, "v"] & Y_{S'}
\end{tikzcd}

Les transformations infinitésimales de type $S'$ de $f$ forment un
\emph{groupe}, contenu dans le produit des groupes des transformations
infinitésimales de type $S'$ de $X$ et $Y$.

Soit $w = (u, v)$ une transformation infinitésimale de type $S'$ de $f$, et
soit $g$ une section de $X/Y$. Considérons alors
\[
  \widetilde{w}(g) = u\, g_{S'}\, v^{-1}
\]
c'est une section de $X_{S'}$ sur $Y_{S'}$, « produit » de $g$, d'où
correspondance : \ill{}

\struck{\ill{}} $X_{S'} \to X$ \ill{} $X \to Y$ \uncertain{soit} $pr_{1}$
\ill{} $Y \to X^{\{S'\}}$ \ill{} $X^{\{S'\}} \to Y^{\{S'\}}$ \ill{}
$Y \to Y^{\{S'\}}$ \ill{}
\note{passage rayé de deux longs traits horizontaux et en partie récrit
au-dessus ; seuls les objets se laissent lire, la syntaxe qui les relie est
perdue}

\[
  \underline{\Gamma}\bigl(X_{Y}^{\{Y_{S'}\}} / Y\bigr)
\]
« produit » de $g$.


[i.e.\ une \ill{} $S$-\ill{} \ill{} \ill{} la relation \ill{}]
\note{trois lignes dans un cadre dessiné à la main, traversé d'un trait
oblique ; c'est l'endroit du lot où un lecteur ayant le feuillet sous les
yeux gagnera le plus}

éléments de $\underline{\Gamma}(X \ill{})$

\emph{Propriétés fonctorielles} \ldots

\page{31}
1) Le groupe des transformations infinitésimales de type $S'/S$ de
$X_{S}/Y_{S}$ \emph{opère} sur \ill{}
\[
  \underline{\Gamma}(X_{S'}/Y_{S'})
\]
des \ill{} de type $S'$ \ill{} sections de $X/Y$, \ill{} de plus
\[
  \underline{\Gamma}(X/Y) \longrightarrow
  \underbrace{\underline{\Gamma}(X_{S'}/Y_{S'})}_{\text{ensemble :
  groupe d'opérateurs}} \simeq
  \underline{\Gamma}\bigl(X_{Y}^{\{Y_{S'}\}} / Y\bigr)
\]
\[
  \operatorname{Aut}^{\{\Lambda\}}_{S}(X'_{S}, Y_{S})
\]
\note{l'accolade et les deux mots « ensemble : groupe d'opérateurs » sont
écrits sous le membre médian ; l'$\operatorname{Aut}$ à exposant
$\{\Lambda\}$ est deux lignes plus bas, isolé}

\textbf{N.B.} $S$, $S'$ \ill{} \ldots $X_{S'} \to Y_{S'}$ ] dans le \ill{}
bijection \ill{} mais seul \ill{} par \ill{} d'opérateurs \ill{} i.e.\ les
$S'$-automorphismes de $Y_{S'}$ et $X_{S'}$ \ldots
\note{le N.B.\ est souligné ; un mot du milieu est biffé et récrit par-dessus,
et l'ensemble n'est lisible que par fragments}

2) Fixons \ill{} \ill{} $w$. \ill{} compatible \ill{}
\[
  w = (u, v), \qquad w' = (u', v')
\]
\ill{} $w$. Un $Y$-morphisme $X \xrightarrow{\;h\;} X'$ \ill{} des \ill{}
\uncertain{transformations}, \ill{} \ill{} rendant commutatif \ill{}

\begin{tikzcd}
  X_{S'} \arrow[r, "h_{S'}"] \arrow[d, "u"'] & X'_{S'} \arrow[d, "u'"] \\
  X_{S'} \arrow[r, "h_{S'}"] & X'_{S'}
\end{tikzcd}

i.e.\ \ill{}

\page{32}
$(u, u')$ \ill{} une \ill{} de type $\Lambda$ de $X \to X'$.

On dit, \ill{} aussi \ill{} pour $g \in \underline{\Gamma}(X/Y)$
\[
  h^{\{Y_{S'}\}}\bigl(\widetilde{w}(g)\bigr) = \widetilde{w}'(h \circ g)
\]
\ill{} un produit : \uncertain{idem} \ill{}

3) Compatibilités avec \ill{}
\[
  w = (u, v), \qquad w' = (u', v')
\]
$X/Y$, $X'/Y$, \ill{} on en déduit \ill{} \ill{} $X \times_{Y} X'$, \ill{}
i.e.\ \ill{} $S'$-automorphismes \ill{}
\[
  \bigl(X \times_{Y} X'\bigr)_{S'} = X_{S'} \times_{Y_{S'}} X'_{S'}
\]
\ill{} de $Y_{S'}$, induisant \ill{} sur $X_{S'}$, $X'_{S'}$ ;
$X_{S'}$, $X'_{S'}$ \ill{} \ill{} structures) \ill{}
\struck{$s_{S'}$} \ill{} $Y_{S'}$ \ill{} $w''$, \ill{}
\[
  \widetilde{w}''\bigl(g \times_{Y} g'\bigr) = \widetilde{w}(g) \times_{Y}
  \widetilde{w}'(g')
\]
l'identification
\[
  \bigl(X \times_{Y} X'\bigr)_{Y}^{\{Y_{S'}\}} \simeq
  X_{Y}^{\{Y_{S'}\}} \times_{Y} X'^{\,\{Y_{S'}\}}_{Y}
\]

\begin{tikzcd}[column sep=small, row sep=small]
  Y_{S'} \arrow[r] & Y_{S'}
\end{tikzcd}

\emph{Comme} \ill{} de composition compatible \ill{} \ill{} \ill{}
\note{le mot souligné qui ouvre le dernier paragraphe est lu « Comme » sans
certitude ; le bas du feuillet est une esquisse, trois ou quatre mots par
ligne}

\section{Exemples, cas particuliers}
\note{titre de sa main, en haut à droite du feuillet 33, chemise par
ailleurs vierge qui couvre les feuillets suivants ; le feuillet ne reçoit pas
de numéro de page}

\page{34}
(i)
\[
  F = a X^{p} + b Y^{p} - \struck{1}
\]
\[
  k[X, Y]/(F) = A \;\simeq\; k[Z]/(Z^{p} - c)\,[X]
\]
\marginal{faire $Z = X + Y$, donc $F = Z^{p} - \ill{}$}
\note{note de sa main en haut à droite, avec au-dessus « $a X^{p} + b Y^{p}$ »
et le mot « non » ; la constante du second membre est écrite $c$ ici et $a$
dans la marge, sans que la page les concilie}

$A$ est un anneau régulier, $k$ est alg.\ fermé dans son corps des fractions,
$A \otimes_{k} \overline{k}$ n'est pas intègre, il est \struck{\ill{}}
isomorphe à
\[
  k[Z]\,[T]/(T^{p})
\]
\note{« $k$ est alg.\ fermé dans son corps des fractions » est cerné d'un
trait ovale de sa main ; dans la dernière formule il écrit $k$ et non
$\overline{k}$, tel quel}

(ii)
\[
  F = Y^{2} - X^{p} + a \qquad (p \neq 2)
\]
\[
  k[X, Y]/(F) = A
\]

$A$ est un anneau régulier, \struck{$k$ est \ill{}} son corps des fractions
est extension séparable de $k$

$A \otimes_{k} \overline{k}$ n'est pas régulier, car le point $Y = 0$,
$X = a^{1/p}$ est singulier sur la courbe $F = 0$
\[
  F = Y^{2} - (X - \lambda)^{p} = 0 \qquad (\lambda = a^{1/p}).
\]
\ill{} donc $k[Y, Z]/(Y^{2} - Z^{p})$, \ill{} l'idéal maximal de l'origine,
où \uncertain{pour} \ill{} : l'origine.
\[
  Y^{2} - Z^{p} \longrightarrow 0
\]
\ill{} polynômes \ill{} \ill{} $=$ \ill{} que $K/k$ \ill{}

Montrons que $A$ est régulier. \ill{} \struck{\ill{} ; \ill{} \ill{}
intégration \ill{}} \ill{}

En effet, la courbe \uncertain{définie} \ill{} $F = Y^{2} - X^{p} + a$
\ill{}

Faisons $X - \lambda = Z$, \ill{}

On a $Y^{2} - Z^{p} \in \mathfrak{m}^{2}$, où $\mathfrak{m}$ \ill{} \ill{}
\ill{} \ill{} peut avoir \ill{}

Cependant, pour $p \neq 2$, \ill{} irréductible. Cela prouve \ill{} une
extension séparable.

\struck{Dans la courbe obtenue \ill{} \ill{} \ill{} un point régulier}
\[
  \Omega^{1}_{k}(A) \simeq A^{2} \Big/ A\Bigl(\frac{\partial F}{\partial X},
  \frac{\partial F}{\partial Y}\Bigr) \simeq A \;\struck{\ill{}}\;
  \simeq A + A/2Y
\]
i.e.\ $\Omega^{1}_{k}(A)$ est \ill{} libre \ill{} $2Y$, i.e.\ les idéaux
\ill{} premiers \ill{} de $k[X, Y]$ \ill{}
\[
  \bigl(Y,\ Y^{2} - X^{p} + a\bigr) = \bigl(Y,\ X^{p} - a\bigr) \;;\qquad
  k[X, Y]/\bigl(Y,\ X^{p} - a\bigr) \simeq k[X]/(X^{p} - a)
\]
\ill{} l'idéal $\mathfrak{m}' = (Y,\ X^{p} - a)$ \ill{} déjà vu \ill{}
c'est \ill{} le

\page{35}
\uncertain{seul} point singulier possible. Or, le \uncertain{carré} de
$\mathfrak{m}'/(F)$ est $\mathfrak{m}''/(F)$
\[
  \mathfrak{m}'' = \bigl(Y^{2},\ Y(X^{p} - a),\ (X^{p} - a)^{2},\
  Y^{2} - X^{p} + a\bigr)
\]
\struck{d'où $Y = \dfrac{X^{p}}{a}$}
\[
  = \bigl((X^{p} - a),\ Y^{2},\ Y(X^{p} - a)\bigr)
\]
\ill{} \uncertain{pour} \ill{}
$\mathfrak{m}'/\mathfrak{m}'^{2} = \mathfrak{m}'/\mathfrak{m}''$ \ill{}
engendré par $Y$ (car $X^{p} - a$ est dans $\mathfrak{m}''$). Donc le point
est \ill{} simple, mais non absolument simple.

Dans ce cas-ci, il vient à \uncertain{vérifier} que $k$ est algébriquement
fermé dans $A$.
\note{le feuillet ne porte que ce paragraphe ; le reste est vierge. « il
vient à vérifier » est lu ainsi, mais les deux mots du milieu sont serrés et
la lecture reste incertaine}

\page{36}
\note{feuillet écrit à l'italienne, sur papier faiblement réglé ; c'est une
feuille de récapitulation, deux colonnes que rien ne relie explicitement}

\emph{Caractérisation des dérivations} \ill{}

[Car \uncertain{relient} : les dérivations \ill{} \ill{} l'adhérence
\uncertain{différentielle} \ill{}]

$\delta f = 0$ : $f$ est \ill{} dans \ill{}

\[
  \struck{d\,\delta f = \delta f \;?}
\]
\[
  \delta f = 0 \iff f \in A[B^{p}]
\]
\[
  \delta^{(\infty)} f = 0 \iff f \in A[B^{p^{k}}]
\]
\note{l'exposant du second critère est lu $(\infty)$ : il est écrit très
petit et pourrait être un $n$ ; ce qui décide en faveur de l'itération, c'est
le $B^{p^{k}}$ du second membre}

$\delta f = 0$, \quad $f \in K^{p}$

Critère de \ill{} \ill{} \ill{} séparation (Chev.\ \uncertain{Weyl})
\[
  \Omega^{1}_{B/A} = 0
\]
\[
  \Omega^{1}_{B/A} = 0
\]
$+$ plat

$\Delta_{B/A}$ est \uncertain{unipotent} \ill{}

$\Delta_{B/A}$ est \ill{}

Critère de \ill{} \ill{} décidable \ill{} localisation (ann.\ \emph{nb.\
$p$-i}) \ill{}

\struck{\ill{}} différentielle \ill{} dérivations de $A \subset B$

\ill{} : $f \in A$

$B/A$ soit une \uncertain{$p$-base}, donc
\[
  \Omega^{1}(B/A)
\]
\note{la colonne de droite est une liste de critères, plusieurs sont
soulignés ; les mots qui les relient sont, pour la plupart, perdus}

\section{Foncteurs différentiels d'ordre $n$}
\note{titre de sa main, en haut à droite du feuillet 38, chemise par ailleurs
vierge ; le feuillet ne reçoit pas de numéro de page}

\page{39}
\emph{Foncteur différentiel d'ordre $\leq n$.}

Soit $S$ un préschéma, $\mathcal{C}$ \struck{une} une \ill{} catégorie
fibrée \uncertain{pleine} de la catégorie \ill{} \ill{} des deux \ill{}
\ill{} \ill{} sur $S'/S$, [la plus \ill{}, \struck{\ill{}} \ill{} \ill{} la
catégorie en question].

Soit $\mathcal{C}'$ une autre catégorie fibrée sur \ill{} $=$ \ill{}
$\mathcal{E}\ell_{S}$ des préschémas sur $S$. [Les \ill{} \uncertain{Faut}
que $\mathcal{C}$ soit \ill{} $\mathcal{C}$ que « verbalement » \ill{}
catégorie de $\mathcal{C}$ \ill{} \struck{plutôt \ill{}} \ill{}
$\mathcal{C}$]

Soit
\[
  F : \mathcal{C} \longrightarrow \Phi^{\mathcal{C}}
\]
\ill{} une \struck{\ill{}} \emph{section continue} \emph{$f$ bornée} \ill{}

$\{X'/S'\}$ \qquad $\{X''/S''\}$ \qquad $\{X'/S'\} \in \Phi_{\mathcal{C}}$

\begin{tikzcd}[column sep=small, row sep=small]
  X' & X'' \arrow[l] \\
  S' \arrow[u] & S'' \arrow[l] \arrow[u] \\
  S \arrow[r] & S' \arrow[r] & S''
\end{tikzcd}
\note{le croquis est de trois lignes ; la ligne du bas, $S \leftarrow S'
\leftarrow S''$, est écrite en dessous et à gauche, rendue ici par une
troisième rangée}

\ill{} $\Phi_{\mathcal{C}}$, \ill{} façon compatible \ill{} \ill{} les
extensions de la base.

\textbf{Ex.} Le foncteur $\underline{P}^{n}_{X'/S'}$ \quad [$\Phi$ le type
\ill{}
\note{la page est un brouillon rapide : les objets portent, les phrases qui
les relient sont pour l'essentiel perdues, et la transcription le dit plutôt
que de les inventer}

\page{40}
\note{le feuillet est traversé de deux longs traits diagonaux, et le tiers
supérieur porte trois couches d'écriture superposées}

\ill{} de \struck{\ill{}} foncteurs \ill{} \ill{} moins : \ill{} \ill{}
foncteurs \ill{} \ill{} les conditions \ill{}

\ill{} foncteur \ill{} \ill{} \ill{}

\[
  F_{\xi} \quad \ill{} \quad \ill{}
\]

On \ill{} groupe \ill{} $(\ldots)_{\xi}$ ] on \ill{} \ill{} foncteurs
$F_{\xi}$, \ill{} \ill{} utilisent \ill{} les foncteurs obtenus \ill{}
\ill{} \struck{\ill{}} \ill{}

\[
  \struck{\ill{} = \ill{}}\qquad \struck{X = \ill{}}
\]

\begin{tikzcd}[column sep=small, row sep=small]
  E^{n} & E^{n}_{S^{d,n}} \arrow[l] & E^{n} \arrow[l] & S \arrow[l] \\
  S & S^{d,n} \arrow[l] & S \arrow[l] &
\end{tikzcd}
\note{les trois colonnes de gauche descendent par des traits verticaux ; une
flèche oblique en pointillé remonte de $S^{d,n}$ vers $E^{n}$, et un signe
d'égalité est placé sous $S^{d,n}$}

\[
  S^{d,n} = \operatorname{Spec}\ \ill{}\, \bigl(A[t_{1}, \ldots, t_{n}] /
  (t_{i} - \ill{})\bigr) \ill{}
\]
\note{la définition de $S^{d,n}$ est écrite sous le diagramme, à demi
couverte par l'un des deux traits diagonaux ; les indices se laissent lire,
l'anneau pas entièrement}

\ill{} \uncertain{automorphismes} \ill{} \ill{} \ill{} façon compatible
\ill{}

\[
  G^{\mathrm{dif}}_{\ill{}}(\ill{})
\]
\ill{} \struck{\ill{}}$_{S}\bigl(\ill{}^{\,d, n}\bigr)$, \ill{}
[\ill{} \; $\varepsilon : S \to S^{d,n}$]

\ill{} \ill{} \ill{}
\[
  E^{n}_{S^{d,n}} \qquad E^{n} \ill{}
\]
\ill{}
\[
  E^{n}_{S^{d,n}} \times_{S^{d,n}} S.
\]
\ill{}
\[
  \bigl\{ E^{n}_{S^{d,n}} \bigr\} = \bigl\{ E^{n} \times_{S} S^{d,n}
  \bigr\}, \quad \ill{}
\]
\ill{} \ill{} quotient \ill{} $E^{n}_{S^{d,n}}$,
\note{la page s'arrête là, sans point final ; c'est le dernier feuillet du
lot}

\end{document}
